\scanpage{49}
\mspage{- 21 -}
% archive sheet number: 48
Tout revient donc à prouver qu'il existe $q$ séries entières absolument sommables par rapport aux $t_i$ et sans terme constant, qui satisfont aux équations
\[
v_k=\frac{M}{(1-\handwritten{\mu}\handwritten{\sum}_{\handwritten{j=1}}^{p}t_j)(1-\handwritten{\mu}\handwritten{\sum}_{\handwritten{i=1}}^{q}v_i)}-M-M\handwritten{\mu}\handwritten{\sum}_{\handwritten{i=1}}^{q}v_i
\qquad (1\handwritten{\leq}k\handwritten{\leq}q)
\tag*{(9)}
\]

Or, on déduit aussitôt \deleted{que}\inserted{\handwritten{de}} (9) que si ce système a une solution,tous les $v_i$ sont égaux à $v_1=v$ sa tisfaisant à
\[
(1-q\handwritten{\mu}v)(1+qM\handwritten{\mu})v=\frac{M}{1-\handwritten{\mu}\handwritten{\sum}_{\handwritten{j=1}}^{p}t_j}-M(1-q\handwritten{\mu}v)
\tag*{(10)}
\]
et réciproquement ; or, la fonction $v(t_1,\ldotp\ldotp,t_p)$ qui s'annule pour les $t_j$ nuls et satisfait à (10) est donnée par
\[
v(t_1,\ldotp\ldotp,t_p)=\frac{1}{2q\handwritten{\mu}(1+qM\handwritten{\mu})}\handwritten{\Biggl[}1-\sqrt{1-4Mq\handwritten{\mu}(1+qM\handwritten{\mu})\left(\frac{1}{1-\handwritten{\mu}\handwritten{\sum}_{\handwritten{j=1}}^{p}t_j}-1\right)}\handwritten{\Biggr]}
\]
Comme la fonction $(1-w)^{\handwritten{\frac{1}{2}}}$ est égale à la somme de la série entière $\handwritten{\sum}_{\handwritten{n=0}}^{\handwritten{\infty}}\binom{\handwritten{\frac{1}{2}}}{n}v^n$ pour $-1\handwritten{<}w\handwritten{<}1$, il résulte du th. de substitution (\handwritten{\S}1,th.2) que $v(t_1,\ldotp\ldotp,t_p)$ est égale à une série entière absolument sommable et sans terme constant dans un polycylindre de centre $(0,\ldots,0)$.

Reste à montrer que si $w_i$ $(1\handwritten{\leq}i\handwritten{\leq}q)$ sont $q$ fonctions \uline{continues} au voisinage de l'origine dans ${K}^{p}$, nulles en ce point et satisfaisant aux relations
\[
w_i=h_{i0}+\handwritten{\sum}_{\handwritten{(n_j)}}h_{i;n_1\ldotp\ldotp n_q}w_1^{n_1}\ldots w_q^{n_q}
\qquad\text{pour }1\handwritten{\leq}i\handwritten{\leq}q
\tag*{(11)}
\]
dans un voisinage de l'origine dans $K^{p}$, on a nécessairement $u_i=w_i$ da ns\handwritten{\unsure{.}} un voisinage de l'origine pour $1\handwritten{\leq}i\handwritten{\leq}q$. Or, compte-tenu de la prop.8 du \handwritten{\S}1, on tire de (5) et (11) que
\[
u_i-w_i=\underline{o}\bigl(\sup_{1\handwritten{\leq}k\handwritten{\leq}q}\handwritten{|}u_k-w_k\handwritten{|}\bigr)
\qquad\text{pour }1\handwritten{\leq}i\handwritten{\leq}q
\]
dans un voisinage de $0$, ce qui entraîne $u_i=w_i$ pour $1\handwritten{\leq}i\handwritten{\leq}q$ dans un voisinage de $0$ (\uline{Fonct.var.réelle}, chap.V,\handwritten{\S}2). Le th.1 est ainsi complètement démontré.
